
This pilot study compared semantic entropy and self-consistency under matched computational budgets across two benchmarks. The key observation is potential benchmark sensitivity: semantic entropy achieved AUROC 0.551 on HaluEval but 0.289 on TruthfulQA, while self-consistency performed near random on both.

\textbf{Summary of Findings:}
\begin{enumerate}
\item Benchmark sensitivity may exist (SE: 0.551 HaluEval vs 0.289 TruthfulQA)
\item Self-consistency with BERTScore performed near random
\item Controlled comparison revealed patterns not visible in isolated evaluations
\end{enumerate}

\textbf{Limitations:}
\begin{itemize}
\item Sample size (N=20) insufficient for statistical conclusions
\item Single model family (Llama-3-8B)
\item Single random seed
\end{itemize}

\textbf{Future Directions:}
\begin{itemize}
\item Full-scale evaluation with $N\geq 200$ samples per dataset
\item Investigation of TruthfulQA labeling methodology alignment
\item Testing NLI-based consistency metrics as alternative to BERTScore
\item Multi-model evaluation (Mistral-7B, other families)
\end{itemize}

This work establishes a matched-budget comparison framework and identifies benchmark-method alignment as a consideration for hallucination detection evaluation. The findings require validation at scale before informing method selection decisions.
